                   IMC 2014, Blagoevgrad, Bulgaria


                                Day 2, August 1, 2014




Problem 1.         For a positive integer x, denote its nth de imal digit by dn(x), i.e. dn (x) ∈
                           ∞
{0, 1, . . . , 9} and x =     dn (x)10n−1 . Suppose that for some sequen e an n=1 , there are
                          P                                                     ∞
                       n=1
only nitely many zeros in the sequen e dn(an ) n=1 . Prove that there are innitely many
                                                           ∞

positive integers that do not o ur in the sequen e (an )∞
                                                        n=1 .
                          (Proposed by Alexander Bolbot, State University, Novosibirsk)
Solution 1.    By the assumption there is some index n0 su h that dn(an ) 6= 1 for n ≥ n0 .
We show that
                           an+1 , an+2 , . . . > 10n for n ≥ n0 .                          (1)
                            ∞
Noti e that in the sum an =    dk (an )10k−1 we have the term dn (an )10n−1 with dn (an ) ≥ 1.
                            P
                               k=1
Therefore, an ≥ 10n−1 . Then for m > n we have am ≥ 10m > 10n. This proves (1).
     From (1) we know that only the rst n elements, a1 , a2 , . . . , an may lie in the interval
[1, 10n ]. Hen e, at least 10n − n integers in this interval do not o ur in the sequen e at
all. As lim(10n − n) = ∞, this shows that there are innitely many numbers that do not
appear among a1 , a2 , . . ..
Solution 2. We will use Cantor's diagonal method to onstru t innitely many positive
integers that do not o ur in the sequen e (an )
     Assume that dn(an ) 6= 0 for n > n0 . Dene the sequen e of digits
                                           (
                                            2 dn (xn ) = 1
                                      gn =
                                            1 dn (xn ) 6= 1.

Hen e gn 6= dn (an ) for every positive integer n. Let
                                    k
                                                          for k = 1, 2, . . . .
                                    X
                             xk =         gn · 10n−1
                                    n=1

    As xk+1 ≥ 10k > xk , the sequen e (xk ) is in reasing and so it ontains innitely many
distin t positive integers. We show that the numbers xn0 , xn0 +1 , xn0 +2 , . . . no not o ur in
the sequen e (an ); in other words, xk 6= an for every pair n ≥ 1 and k ≥ n0 of integers.
    Indeed, if k ≥ n then dn (xk ) = gn 6= dn(an ), so xk 6= an .
    If n > k ≥ n0 then dn (xk ) = 0 6= dn (an ), so xk 6= an .


Problem 2.            Let A = (aij )ni,j=1 be a symmetri n × n matrix with real entries, and let
λ1 , λ2 , . . . , λn denote its eigenvalues. Show that
                                     X                      X
                                              aii ajj ≥             λi λj ,
                                    1≤i<j≤n               1≤i<j≤n

and determine all matri es for whi h equality holds.
                                     (Proposed by Martin Niepel, Comenius University, Bratislava)
Solution.   Eigenvalues of a real symmetri matrix are real, hen e the inequality makes
sense. Similarly, for Hermitian matri es diagonal entries as well as eigenvalues have to be
real.
    Sin e the tra e of a matrix is the sum of its eigenvalues, for A we have
                                                      n
                                                      X               n
                                                                      X
                                                            aii =           λi ,
                                                      i=1             i=1

and onsequently
                                  n
                                  X                X                  n
                                                                      X               X
                                        a2ii + 2          aii ajj =         λ2i + 2          λi λj .
                                  i=1               i<j               i=1              i<j

   Therefore our inequality is equivalent to
                                                      n
                                                      X               n
                                                                      X
                                                            a2ii ≤          λ2i .
                                                      i=1             i=1

   Matrix A , whi h is equal to A A (or A A in Hermitian ase), has eigenvalues λ21 , λ22 , . . . , λ2n.
               2                               T              ∗

On the other hand, the tra e of AT A gives the square of the Frobenius norm of A, so we
have                n          n                               n
                                                   |aij |2 = tr(AT A) = tr(A2 ) =
                            X              X                                                   X
                                  a2ii ≤                                                               λ2i .
                            i=1            i,j=1                                                i=1

   The inequality follows, and it is lear that the equality holds for diagonal matri es
only.
Remark.    Same statement is true for Hermitian matri es.


                           sin x
Problem 3.         Let f (x) =   , for x > 0, and let n be a positive integer. Prove that
                             x
             1
f (n) (x) <     , where f (n) denotes the nth derivative of f .
            n+1
                                    (Proposed by Alexander Bolbot, State University, Novosibirsk)
Solution 1.       Putting f (0) = 1 we an assume that the fun tion is analyti in R. Let
g(x) = x n+1       n     1
               (f (x) − n+1 ). Then g(0) = 0 and
                                                                              
                                                                       1
                       g (x) = (n + 1)x f (n) (x) −
                        ′                      n
                                                                                    + xn+1 f (n+1) (x) =
                                                                      n+1
                                                             
= xn (n + 1)f (n) (x) + xf (n+1) (x) − 1 = xn (xf (x))(n+1 − 1 = xn (sin(n+1) (x) − 1) ≤ 0.
                                        

Hen e g(x) ≤ 0 for x > 0. Taking into a ount that g ′ (x) < 0 for 0 < x < π2 we obtain
the desired (stri t) inequality for x > 0.
Solution 2.

                   (n)           Z 1                       Z 1                            Z 1
                            dn                                     ∂n
       
           sin x
                          = n             − cos(xt)dt =               (− cos(xt)) dt =           tn gn (xt)dt
             x             dx         0                      0    ∂xn                       0

where the fun tion gn (u) an be ± sin u or ± cos u, depending on n. We only need that
|gn | ≤ 1 and equality o urs at nitely many points. So,
                                     (n)       Z 1                    Z 1
                              sin x                    n                                 1
                                             ≤         t gn (xt) dt <         tn dt =       .
                                x                 0                      0              n+1



Problem 4.     We say that a subset of Rn is k-almost ontained by a hyperplane if there
are less than k points in that set whi h do not belong to the hyperplane. We all a nite
set of points k-generi if there is no hyperplane that k-almost ontains the set. For ea h
pair of positive integers k and n, nd the minimal number d(k, n) su h that every nite
k -generi set in Rn ontains a k -generi subset with at most d(k, n) elements.
(Proposed by Sha har Carmeli, Weizmann Inst. and Lev Radzivilovsky, Tel Aviv Univ.)
                                   (
                                    k · n k, n > 1
Solution. The answer is: d(k, n) =
                                    k + n otherwise
    Throughout the solution, we shall often say that a hyperplanes skips a point to signify
that the plane does not ontain that point.
    For n = 1 the laim is obvious.
    For k = 1 we have an arbitrary nite set of points in Rn su h that neither hyperplane
 ontains it entirely. We an build a subset of n + 1 points step by step: on ea h step we
add a point, not ontained in the minimal plane spanned by the previous points. Thus
any 1-generi set ontains a non-degenerate simplex of n + 1 points, and obviously a
non-degenerate simplex of n + 1 points annot be redu ed without loosing 1-generality.
    In the ase k, n > 1 we shall give an example of k · n points. On ea h of the Cartesian
axes hoose k distin t points, dierent from the origin. Let's show that this set is k-
generi . There are two types of planes: ontaining the origin and skipping it. If a plane
 ontains the origin, it either ontains all the hose points of a axis or skips all of them.
Sin e no plane ontains all axes, it skips the k hosen points on one of the axes. If a plane
skips the origin, it it ontains at most one point of ea h axis. Therefore it skips at least
n(k − 1) points. It remains to verify a simple inequality n(k − 1) ≥ k whi h is equivalent
to (n − 1)(k − 1) ≥ 1 whi h holds for n, k > 1.
    The example we have shown is minimal by in lusion: if any point is removed, say a
point from axis i, then the hyperplane xi = 0 skips only k − 1 points, and our set stops
being k-generi . Hen e d(k, n) ≥ kn.
    It remains to prove that Hen e d(k, n) ≥ kn for k, n > 1, meaning: for ea h k-generi
nite set of points, it is possible to hoose a k-generi subset of at most kn points. Let
us all a subset of points minimal if by taking out any point, we loose k-generality.
It su es to prove that any minimal k-generi subset in Rn has at most kn points. A
hyperplane will be alled ample if it skips pre isely k points. A point annot be removed
from a k-generi set, if and only if it is skipped by an ample hyperplane. Thus, in a
minimal set ea h point is skipped by an ample hyperplane.
    Organise the following pro ess: on ea h step we hoose an ample hyperplane, and paint
blue all the points whi h are skipped by it. Ea h time we hoose an ample hyperplane,
whi h skips one of the unpainted points. The unpainted points at ea h step (after the
beginning) is the interse tion of all hosen hyperplanes. The interse tion set of hosen
hyperplanes is redu ed with ea h step (sin e at least one point is being painted on ea h
step).
    Noti e, that on ea h step we paint at most k points. So if we start with a minimal
set of more then nk points, we an hoose n planes and still have at least one unpainted
points. The interse tion of the hosen planes is a point (sin e on ea h step the dimension
of the interse tion plane was redu ed), so there are at most nk + 1 points in the set. The
last unpainted point will be denoted by O . The last unpainted line (whi h was formed on
the step before the last) will be denoted by ℓ1 .
    This line is an interse tion of all the hosen hyperplanes ex ept the last one. If we
have more than nk points, then ℓ1 ontains exa tly k + 1 points from the set, one of whi h
is O .
    We ould have exe uted the same pro ess with hoosing the same hyperplanes, but in
dierent order. Anyway, at ea h step we would paint at most k points, and after n steps
only O would remain unpainted; so it was pre isely k points on ea h step. On step before
the last, we might get a dierent line, whi h is interse tion of all planes ex ept the last
one. The lines obtained in this way will be denoted ℓ1 , ℓ2 , ..., ℓn , and ea h ontains exa tly
k points ex ept O . Sin e we have O and k points on n lines, that is the entire set. Noti e
that the ve tors spanning these lines are linearly independent (sin e for ea h line we have
a hyperplane ontaining all the other lines ex ept that line). So by removing O we obtain
the example that we've des ribed already, whi h is k-generi .
Remark.     From the proof we see, that the example is unique.


Problem 5.            For every positive integer n, denote by Dn the number of permutations
(x1 , . . . , xn ) of (1, 2, . . . , n) su h that xj 6= j for every 1 ≤ j ≤ n. For 1 ≤ k ≤ n2 , denote
by ∆(n, k) the number of permutations (x1 , . . . , xn ) of (1, 2, . . . , n) su h that xi = k + i
for every 1 ≤ i ≤ k and xj 6= j for every 1 ≤ j ≤ n. Prove that
                                              k−1 
                                                    k − 1 D(n+1)−(k+i)
                                              X          
                                 ∆(n, k) =                                   .
                                              i=0
                                                       i      n − (k + i)

         (Proposed by Combinatori s; Ferdowsi University of Mashhad, Iran; Mirzavaziri)
Solution.     Let ar ∈ {i1 , . . . , ik } ∩ {a1 , . . . , ak }. Thus ar = is for some s 6= r. Now there
are two ases:
    Case 1. as ∈ {i1 , . . . , ik }. Let as = it . In this ase a derangement x = (x1 , . . . , xn )
satises the ondition xij = aj if and only if the derangement x′ = (x′1 , . . . , x′it −1 , x′it +1 , x′n )
of the set [n] \ {it } satises the ondition x′ij = a′j for all j 6= t, where a′j = aj for j 6= s
and a′s = at . This provides a one to one orresponden e between the derangements
x = (x1 , . . . , xn ) of [n] with xij = aj for the given sets {i1 , . . . , ik } and {a1 , . . . , ak } with
ℓ elements in their interse tions, and the derangements x′ = (x′1 , . . . , x′it −1 , x′it +1 , x′n ) of
[n] \ {it } with xij = a′j for the given sets {i1 , . . . , ik } \ {it } and {a′1 , . . . , a′k } \ {a′t } with
ℓ − 1 elements in their interse tions.
    Case 2. as ∈/ {i1 , . . . , ik }. In this ase a derangement x = (x1 , . . . , xn ) satises the
 ondition xij = aj if and only if the derangement x′ = (x′1 , . . . , x′as −1 , x′as +1 , x′n ) of the
set [n] \ {as} satises the ondition x′ij = aj for all j 6= s. This provides a one to one
 orresponden e between the derangements x = (x1 , . . . , xn ) of [n] with xij = aj for the
given sets {i1 , . . . , ik } and {a1 , . . . , ak } with ℓ elements in their interse tions, and the
derangements x′ = (x′1 , . . . , x′as −1 , x′as +1 , x′n ) of [n] \ {as } with xij = aj for the given sets
{i1 , . . . , ik } \ {is } and {a1 , . . . , ak } \ {as } with ℓ − 1 elements in their interse tions.
     These onsiderations show that ∆(n, k, ℓ) = ∆(n − 1, k − 1, ℓ − 1). Iterating this
argument we have
                                  ∆(n, k, ℓ) = ∆(n − ℓ, k − ℓ, 0).
We an therefore assume that ℓ = 0. We thus evaluate ∆(n, k, 0), where 2k ⩽ n. For
k = 0, we obviously have ∆(n, 0, 0) = Dn . For k ⩾ 1, we laim that

                      ∆(n, k, 0) = ∆(n − 1, k − 1, 0) + ∆(n − 2, k − 1, 0).

For a derangement x = (x1 , . . . , xn ) satisfying xij = aj there are two ases: xa1 = i1 or
xa1 6= i1 .
    If the rst ase o urs then we have to evaluate the number of derangements of the
set [n] \ {i1 , a1 } for the given sets {i2 , . . . , ik } and {a2 , . . . , ak } with 0 elements in their
interse tions. The number is equal to ∆(n − 2, k − 1, 0).
    If the se ond ase o urs then we have to evaluate the number of derangements of
the set [n] \ {a1 } for the given sets {i2 , . . . , ik } and {a2 , . . . , ak } with 0 elements in their
interse tions. The number is equal to ∆(n − 1, k − 1, 0).
    We now use indu tion on k to show that
                                    k−1 
                                          k − 1 D(n+1)−(k+i)
                                    X          
                     ∆(n, k, 0) =                                 ,   2 ⩽ 2k ⩽ n.
                                     i=0
                                             i      n − (k + i)

For k = 1 we have
                                                                                       Dn
            ∆(n, 1, 0) = ∆(n − 1, 0, 0) + ∆(n − 2, 0, 0) = Dn−1 + Dn−2 =                  .
                                                                                      n−1
Now let the result be true for k − 1. We an write
 ∆(n, k, 0) = ∆(n − 1, k − 1, 0) + ∆(n − 2, k − 1, 0)
              k−2                                       k−2 
                                 Dn−(k−1+i)                              D(n−1)−(k−1+i)
                                                                  
              X     k−2                                  X    k−2
            =                                         +
              i=0
                      i    (n − 1) − (k − 1 + i) i=0            i     (n − 2) − (k − 1 + i)
              k−2                             k−1 
                    k − 2 D(n+1)−(k+i) X k − 2                    Dn−(k+i−1)
              X                                           
            =                               +
              i=0
                      i    n − (k + i)         i=1
                                                     i − 1 (n − 1) − (k + i − 1)
                          k−2 
              D(n+1)−k X k − 2 D(n+1)−(k+i)
                                       
            =           +
                n−k       i=1
                                   i        n − (k + i)
                                 k−2 
                D(n+1)−(2k−1) X k − 2 D(n+1)−(k+i)
                                                
              +               +
                n − (2k − 1)     i=1
                                         i − 1 n − (k + i)
                          k−2 
              D(n+1)−k X  k − 2                  k − 2  D(n+1)−(k+i) D(n+1)−(2k−1)
                                                     
            =           +                    +                         +
                n−k       i=1
                                     i            i−1      n − (k + i)    n − (2k − 1)
                          k−2 
              D(n+1)−k X k − 1 D(n+1)−(k+i) D(n+1)−(2k−1)
                                       
            =           +                                 +
                n−k       i=1
                                   i        n −  (k + i)    n − (2k − 1)
              k−1 
                    k − 1 D(n+1)−(k+i)
              X          
            =                              .
              i=0
                      i    n  − (k  +  i)

Remark.    As a orollary of the above problem, we an solve the rst problem. Let n = 2k,
ij = j and aj = k + j for j = 1, . . . , k. Then a derangement x = (x1 , . . . , xn ) satises the
 ondition xij = aj if and only if x′ = (xk+1 , . . . , xn ) is a permutation of [k]. The number of su h
permutations x′ is k!. Thus i=0            k−i = k!.
                             Pk−1 k−1 Dk+1−i
                                       i
